\scanpage{310}

\mspage{49}
% archive sheet number: 309

Cela implique que non seulement $\operatorname{Ob}\underline{\mathrm{Ty}}=\mathrm{Ty}$, mais aussi $\operatorname{Fl}(\underline{\mathrm{Ty}})$, et plus généralement tous les ensembles \deleted{\unsure{de|tels}} de la forme $\operatorname{Hom}_{\mathrm{Cat}}(X,\underline{\mathrm{Ty}})$, \inserted{sont des monoïdes, et ceci fonctoriellement en $X$ ;} par exemple les composantes $\operatorname{Hom}_{\mathrm{Cat}}(\Delta_i,\underline{\mathrm{Ty}})$ de l'ens. semi-simplicial ``nerf'' de $\underline{\mathrm{Ty}}$
\[
\operatorname{Nerf}(\underline{\mathrm{Ty}})_i=\operatorname{Hom}_{\mathrm{Cat}}(\Delta_{\unsure{i}},\underline{\mathrm{Ty}}),
\]
sont des monoïdes, et les applications \unsure{sesimpl.|ssimpl.} entre elles sont des hom. de monoïdes : le \unsure{Nerf} \deleted{est} de $\underline{\mathrm{Ty}}$ est un \uline{monoïde} \uline{semi-simplicial}.

On va décrire ces monoïdes $\operatorname{Nerf}_i(\underline{\mathrm{Ty}})$, et constater que ce sont \inserted{tous} des monoïdes \uline{libres}, \deleted{\unsure{et}} \deleted{\unsure{explicit.}} tout comme $\operatorname{Nerf}_0(\underline{\mathrm{Ty}})=\operatorname{Ob}\underline{\mathrm{Ty}}=\mathrm{Ty}$. Les applications \unsure{sspl.} entre ces monoïdes sont donc connues quand on les connaît sur les \uline{générateurs} de ces monoïdes libres, (\deleted{\unsure{c. à d. sur les}}\inserted{lesquels générateurs} sont aussi les éléments indécomposables, donc définis intrinsèquement en termes de \unsure{sa|la} structure multiplicative).

Pour $\operatorname{Nerf}_0(\underline{\mathrm{Ty}})=\mathrm{Ty}$, c'est une chose vue, les générateurs sont
\[
\Delta_+,\Delta_-,
\]
correspondant respectivement aux intervalles orientés
\diagram{\begin{tikzcd}[ampersand replacement=\&,row sep=0pt,column sep=small]
\Delta_+: \& 0 \arrow[r] \& 1\\
\Delta_-: \& 0 \& 1 \arrow[l]
\end{tikzcd}.}
